\subsection{Аналіз швидкості збіжності інтерполяційного процесу}
\paragraph{теорема Вейєрштрасса.} Фундаментальною основою для використання поліномів у задачах наближення функцій є класичне твердження, яке гарантує принципову можливість такого наближення. Воно стверджує, що будь-яка неперервна функція на обмеженому інтервалі може бути наближена поліномами з довільною наперед заданою точністю. Цей факт формалізується у вигляді відомої теореми\cite{trefethen2013approximation}.

\begin{statement}[Апроксимаційна теорема Вейєрштрасса]
	\label{st:weierstrass}
	Нехай $f$~--- неперервна функція на відрізку $[-1, 1]$, і нехай задано довільне як завгодно мале додатне число $\varepsilon > 0$. Тоді існує такий поліном $p$, для якого виконується умова:
	\begin{equation}\label{eq:weierstrass_condition}
		\|f - p\| < \varepsilon.
	\end{equation}
\end{statement}

Важливим наслідком теореми Вейєрштрасса є те, що вона встановлює принципову можливість поліноміального наближення навіть для вкрай негладких функцій. До цієї категорії належать, наприклад, такі функції як $x \sin(x^{-1})$ або навіть $\sin(x^{-1})\sin(1/\sin(x^{-1}))$ ~\cite{trefethen2013approximation}.

Центральною догмою теорії наближень можна вважати наступний фундаментальний принцип: чим більш гладкою є функція, тим швидше збігаються її апроксимації при $n \to \infty$. 

Нехай функція $f$ є $\nu$ разів диференційовною на відрізку $[-1, 1]$ (можливо, зі скачками у $\nu$-й похідній). При дослідженні збіжності її чебишовських інтерполянтів при $n \to \infty$ у рівномірній $\infty$-нормі, швидкість збіжності зазвичай становить $\mathcal{O}(n^{-\nu})$. Показовим прикладом є функція $f(x) = |x|$, яка є один раз диференційовною зі скачком першої похідної в точці $x = 0$. Її крива збіжності ідеально збігається зі швидкістю $n^{-1}$, причому максимальна похибка інтерполянта $p_n$ досягається в нулі і для непарних $n$ точно дорівнює $p(0) = 1/n$.

При спробі знайти суворе теоретичне обґрунтування оцінки $\mathcal{O}(n^{-\nu})$ в класичній літературі виникають труднощі. Найвідоміші теореми Джексона вимагають \textit{строгої неперервності} похідних. Проте для наведеного прикладу $f(x) = |x|$ перша похідна не є неперервною. Якби ми спиралися лише на класичні теореми, нам довелося б задовольнитися нульовою похідною, що гарантувало б лише збіжність $\mathcal{O}(n^0)$, а не фактично спостережувану $\mathcal{O}(n^{-1})$.

Щоб отримати точну оцінку, необхідно спиратися на властивість, яку має більшість функцій у прикладних задачах — \textit{обмежену варіацію} (bounded variation). Функція (незалежно від того, неперервна вона чи ні) має обмежену варіацію, якщо її повна варіація є скінченною. Повна варіація визначається як 1-норма її похідної: $V = \|f'\|_1$. Саме скінченність цієї величини дозволяє інтерполянтам досягати високої швидкості збіжності.

Для формування строгої математичної бази почнемо з оцінки самих коефіцієнтів Чебишева. Умова \textit{абсолютної неперервності}, що використовується нижче, є стандартною: вона означає, що функція дорівнює інтегралу від своєї похідної, яка існує майже скрізь і є інтегровною за Лебегом.

\begin{statement}[Коефіцієнти Чебишева для диференційовних функцій]
	\label{st:cheb_coeffs_diff}
	Для цілого $\nu \ge 0$, нехай функція $f$ та її похідні до порядку $f^{(\nu-1)}$ є абсолютно неперервними на $[-1, 1]$, і нехай $\nu$-та похідна $f^{(\nu)}$ має обмежену повну варіацію $V$. Тоді для $k \ge \nu + 1$ коефіцієнти Чебишева функції $f$ задовольняють нерівність:
	\begin{equation}\label{eq:cheb_coeff_bound_full}
		|a_k| \le \frac{2V}{\pi k(k-1)\cdots(k-\nu)} \le \frac{2V}{\pi(k-\nu)^{\nu+1}}.
	\end{equation}
\end{statement}

\begin{proof}
	Використовуючи підстановку $x = \frac{1}{2}(z + z^{-1})$, де $z$ лежить на одиничному колі у комплексній площині, коефіцієнт визначається як:
	\begin{equation*}
		a_k = \frac{1}{\pi i} \int_{|z|=1} f\left(\frac{1}{2}(z + z^{-1})\right) z^{k-1} dz.
	\end{equation*}
	Інтегрування частинами по $z$ перетворює цей вираз на:
	\begin{equation}\label{eq:ak_int_parts}
		a_k = \frac{-1}{\pi i} \int_{|z|=1} f'\left(\frac{1}{2}(z + z^{-1})\right) \frac{z^k}{k} \frac{dx}{dz} dz.
	\end{equation}
	Множник $dx/dz$ з'являється через те, що $f'$ позначає похідну по $x$, а не по $z$. 
	
	Розглянемо випадок $\nu = 0$, де припускається лише те, що $f$ має обмежену варіацію $V = \|f'\|_1$. Інтеграл по верхній половині одиничного кола еквівалентний інтегралу по $x$; інтеграл по нижній половині дає ще один такий самий інтеграл. Об'єднуючи їх, отримуємо:
	\begin{equation*}
		a_k = \frac{1}{\pi i} \int_{-1}^1 f'(x) \frac{z^k - z^{-k}}{k} dx = \frac{2}{\pi} \int_{-1}^1 f'(x) \Im\left(\frac{z^k}{k}\right) dx.
	\end{equation*}
	Оскільки $|z^k/k| \le 1/k$ для $x \in [-1, 1]$, а $V = \|f'\|_1$, з цього безпосередньо випливає $|a_k| \le 2V/\pi k$.
	
	Якщо $\nu > 0$, замінимо $dx/dz$ на $\frac{1}{2}(1 - z^{-2})$ у рівнянні~\eqref{eq:ak_int_parts}, отримуючи:
	\begin{equation*}
		a_k = -\frac{1}{\pi i} \int_{|z|=1} f'\left(\frac{1}{2}(z + z^{-1})\right) \left[ \frac{z^k}{2k} - \frac{z^{k-2}}{2k} \right] dz.
	\end{equation*}
	Повторне інтегрування частинами по $z$ дає:
	\begin{equation*}
		a_k = \frac{1}{\pi i} \int_{|z|=1} f''\left(\frac{1}{2}(z + z^{-1})\right) \left[ \frac{z^{k+1}}{2k(k+1)} - \frac{z^{k-1}}{2k(k-1)} \right] \frac{dx}{dz} dz.
	\end{equation*}
	Для $\nu = 1$ (де $f'$ має варіацію $V = \|f''\|_1$), цей інтеграл знову зводиться до інтеграла по $x$:
	\begin{equation*}
		a_k = \frac{-2}{\pi} \int_{-1}^1 f''(x) \Im\left[ \frac{z^{k+1}}{2k(k+1)} - \frac{z^{k-1}}{2k(k-1)} \right] dx.
	\end{equation*}
	Оскільки вираз у квадратних дужках обмежений величиною $1/k(k-1)$, маємо $|a_k| \le 2V/\pi k(k-1)$. Для вищих похідних ($\nu > 1$) процес інтегрування частинами продовжується аналогічно загалом $\nu + 1$ разів, що на кожному кроці додає відповідні множники у знаменник і доводить загальну теорему.
\end{proof}

З цієї теореми можна вивести прямі наслідки щодо точності самих чебишовських інтерполянтів.

\begin{statement}[Збіжність для диференційовних функцій]
	\label{st:convergence_differentiable}
	Якщо функція $f$ задовольняє умови попередньої теореми, де $V$ позначає повну варіацію $f^{(\nu)}$ для деякого $\nu \ge 1$, то для будь-якого $n > \nu$ її чебишовські проєкції $f_n$ задовольняють оцінку:
	\begin{equation}\label{eq:cheb_proj_bound}
		\|f - f_n\| \le \frac{2V}{\pi\nu(n-\nu)^\nu},
	\end{equation}
	а її чебишовські інтерполянти $p_n$ задовольняють оцінку:
	\begin{equation}\label{eq:cheb_interp_bound}
		\|f - p_n\| \le \frac{4V}{\pi\nu(n-\nu)^\nu}.
	\end{equation}
\end{statement}

\begin{proof}
	Для доведення оцінки~\eqref{eq:cheb_proj_bound} застосуємо нерівність для коефіцієнтів із твердження~\ref{st:cheb_coeffs_diff} до мажорантного ряду похибки проєкції:
	\begin{equation}\label{eq:proof_proj_sum}
		\| f - f_n \| \le \sum_{k=n+1}^{\infty} |a_k| \le \frac{2V}{\pi} \sum_{k=n+1}^{\infty} (k-\nu)^{-\nu-1}.
	\end{equation}
	Цю суму, у свою чергу, можна обмежити зверху інтегралом:
	\begin{equation}\label{eq:proof_integral_bound}
		\int_n^\infty (s-\nu)^{-\nu-1} ds = \frac{1}{\nu(n-\nu)^\nu}.
	\end{equation}
	Для доведення оцінки~\eqref{eq:cheb_interp_bound} використовується аналогічний розклад для похибки інтерполяції, який замість поодиноких коефіцієнтів містить додавання подвоєних значень $2|a_k|$ замість $|a_k|$, що безпосередньо призводить до появи множника $4$ у чисельнику фінального виразу.
\end{proof}

Коротко кажучи: наявність $\nu$-ї похідної з обмеженою варіацією означає збіжність з алгебраїчною швидкістю $\mathcal{O}(n^{-\nu})$. 

\paragraph{Практичне значення для обробки зображень.}
З точки зору цифрової обробки сигналів, наведений математичний аналіз має фундаментальне значення для обґрунтування вибору методу масштабування. Цифрове зображення за своєю природою не є нескінченно гладкою функцією. Навпаки, реальні фотографії містять безліч різких переходів інтенсивності (контури об'єктів, межі тіней, висококонтрастні текстури), які математично еквівалентні функціям із низьким ступенем диференційовності або наявністю скачків у похідних. 

Якби для інтерполяції таких даних використовувалася класична поліноміальна модель на рівномірній сітці, наявність різких контурів провокувала б глобальне розкачування полінома. Це неминуче призвело б до катастрофічного накопичення похибки (феномен Рунге) та появи сильних візуальних спотворень у вигляді паразитних хвиль навколо контурів (артефакти дзвону, ringing artifacts).

Теореми~\ref{st:cheb_coeffs_diff} та~\ref{st:convergence_differentiable} строго доводять, що використання чебишовських вузлів (які лежать в основі розробленого барицентричного алгоритму) гарантує стабільну алгебраїчну збіжність $\mathcal{O}(n^{-\nu})$ навіть для функцій, що мають лише обмежену варіацію, а не ідеальну гладкість. Це означає, що при масштабуванні зображення розробленим методом похибка локалізується, поліном не зазнає експоненціального розкачування на краях, а різкі контури об'єктів залишаються стабільними. 